\documentclass[11pt]{article}

\usepackage[margin=1in]{geometry}
\usepackage{amsmath}
\usepackage{amssymb}
\usepackage{amsthm}
\usepackage{mathrsfs}
\usepackage{mathtools}
\usepackage{tikz}
\usepackage[colorlinks=true, linkcolor=blue, pdfusetitle]{hyperref}

\renewcommand{\qedsymbol}{$\blacksquare$}

\title{\textbf{Advanced Practice Problem Set: Existence Proofs, Logic, and Quantitative Modeling}}
\author{Tashfeen Omran}
\date{\today}

\begin{document}

\maketitle

\tableofcontents
\clearpage

\section{Concept Checklist \& Question Mapping}
\label{sec:checklist_start}
This problem set rigorously tests the abstract concepts, logical structures, and proof techniques extracted from the foundational text on Existence Proofs and Disproving Existence Statements. 

\begin{itemize}
    \item \textbf{C1: Constructive Existence (Algebraic/Rationality)} -- Problems 1, 13, 25
    \item \textbf{C2: Constructive Existence (Primes/Divisibility/Factorials)} -- Problems 2, 14, 26
    \item \textbf{C3: Non-Constructive Existence (Intermediate Value Theorem)} -- Problems 3, 8, 15, 27
    \item \textbf{C4: Combinatorial Existence (Pigeonhole Principle)} -- Problems 4, 16, 28, 30
    \item \textbf{C5: Disproving Existence (Algebra \& Sum of Squares)} -- Problems 5, 17, 29
    \item \textbf{C6: Disproving Existence (Parity \& Consecutive Integers)} -- Problems 6, 18
    \item \textbf{C7: Disproving Existence (Modular Arithmetic \& Finite Fields)} -- Problems 7, 19
    \item \textbf{C8: Uniqueness Proofs ($\exists!$)} -- Problems 8, 10, 20
    \item \textbf{C9: Set-Theoretic \& Logical Equivalences (De Morgan's)} -- Problems 9, 21
    \item \textbf{C10: Abstract Real-World Modeling (Quant Finance/Risk-Neutral Measures)} -- Problems 10, 22
    \item \textbf{C11: Formal Limit Existence ($\epsilon-\delta$ Nested Quantifiers)} -- Problems 11, 23
    \item \textbf{C12: Advanced Diagnostic Testing (Finding the Flaw)} -- Problems 12, 24
\end{itemize}

\clearpage

\section{Problem Set}

\begin{enumerate}
    \item \textbf{[C1]} Prove constructively that there exist irrational numbers $a$ and $b$ such that $a^b = 2$. 
    \item \textbf{[C2]} Prove that there exist six distinct positive integers such that each integer divides the sum of the remaining five integers.
    \item \textbf{[C3]} Prove that the equation $x^5 - 3x^2 + 1 = 0$ has at least one real root in the interval $[0, 1]$.
    \item \textbf{[C4]} Prove that if $5$ points are placed inside a square of side length $2$, at least two points are within a distance of $\sqrt{2}$ from each other.
    \item \textbf{[C5]} Disprove the following statement: There exists a real number $x$ such that $x^4 + 2x^2 + 2 = 2x$.
    \item \textbf{[C6]} Disprove the following statement: There exists an integer $n$ such that $n^6 + n^3 + 1$ is an even integer.
    \item \textbf{[C7]} Disprove the following statement: There exist odd integers $a, b, c$ such that $a^2 + b^2 + c^2 = 8k + 7$ for some integer $k$.
    \item \textbf{[C8, C3]} Prove that there exists a unique real number $x$ such that $e^x + x = 0$.
    \item \textbf{[C9]} Let $S$ be a domain. Establish the logical equivalence of the following statements by manipulating quantifiers: $\sim(\exists x \in S, P(x) \implies Q(x)) \equiv \forall x \in S, P(x) \land \sim Q(x)$.
    \item \textbf{[C10, C8]} In a one-period binomial market model, let the initial stock price be $S_0 = 50$. At time $t=1$, the stock can move to $S_1(\omega_1) = 60$ or $S_1(\omega_2) = 40$. The risk-free interest rate is $r = 0.1$. Prove that there exists a unique risk-neutral probability measure $\mathbb{Q}$.
    \item \textbf{[C11]} Prove the existence of the following limit using the formal $\epsilon-\delta$ definition: $\lim_{x \to 4} (3x - 2) = 10$.
    \item \textbf{[C12]} Find the logical flaw in the following proof: \textit{Statement: There exists a real number $x$ such that $x^2 + 1 = 0$. Proof: Let $x = i$. Then $i^2 + 1 = -1 + 1 = 0$. Thus, such an $x$ exists.}
    \item \textbf{[C1]} Prove that there exist a rational number $a$ and an irrational number $b$ such that $a^b$ is irrational.
    \item \textbf{[C2]} Prove that there exists a prime number $p$ such that $p, p+2$, and $p+6$ are all prime numbers.
    \item \textbf{[C3]} Prove that the equation $\cos(x) = x^2$ has at least two distinct real roots.
    \item \textbf{[C4]} Prove that among any $n+1$ integers, there exist two distinct integers whose difference is a multiple of $n$.
    \item \textbf{[C5]} Disprove the following statement: There exist real numbers $x$ and $y$ such that $x^2 + y^2 - 4x - 6y + 14 = 0$.
    \item \textbf{[C6]} Disprove the following statement: There exist integers $a$ and $b$ such that $a^2 - b^2 = 2$.
    \item \textbf{[C7]} Disprove the following statement: There exist integers $x$ and $y$ such that $x^2 - 3y^2 = 8$.
    \item \textbf{[C8]} Prove that there exists a unique set $A$ such that for all sets $B$, $A \cap B = A$.
    \item \textbf{[C9]} Prove that the statement $\exists x \in \mathbb{R}, \forall y \in \mathbb{R}, x + y = 0$ is false.
    \item \textbf{[C10]} In a one-period binomial model, prove that if the up-factor $u$ equals $1+r$ (where $r$ is the risk-free rate), there does not exist an equivalent risk-neutral measure $\mathbb{Q}$ (where $0 < \mathbb{Q}(\omega_1) < 1$).
    \item \textbf{[C11]} Prove the existence of the following limit using the formal $\epsilon-\delta$ definition: $\lim_{x \to 2} x^2 = 4$.
    \item \textbf{[C12]} Find the logical flaw in the following proof: \textit{Statement: $\exists x \in \mathbb{R}$ such that $x > x+1$. Proof: Assume for contradiction that $\forall x \in \mathbb{R}, x \le x+1$. Subtracting $x$ from both sides yields $0 \le 1$, which is a true statement. Since the negation leads to a true statement, the original statement must be true.}
    \item \textbf{[C1]} Prove that there exist infinitely many rational numbers in the open interval $(0, 1)$.
    \item \textbf{[C2]} Prove that there exists a positive integer $n$ such that $n! + 1$ is a perfect square.
    \item \textbf{[C3]} Prove Brouwer's Fixed Point Theorem for one dimension: Any continuous function $f: [a, b] \to [a, b]$ has at least one fixed point $c \in [a, b]$ such that $f(c) = c$.
    \item \textbf{[C4]} Prove that if $101$ items are placed into $10$ containers, there exists at least one container that holds $11$ or more items.
    \item \textbf{[C5]} Disprove the following statement: There exist strictly positive integers $x$ and $y$ such that $x^2 - y^2 = 1$.
    \item \textbf{[C4]} Let $S$ be a set of $5$ distinct positive integers. Prove that there exist two distinct non-empty subsets $A$ and $B$ of $S$ such that the sum of the elements in $A$ is congruent to the sum of the elements in $B$ modulo $30$.
\end{enumerate}

\clearpage

\section{Detailed Solutions}

\begin{enumerate}
    % Solution 1
    \item \textbf{Proof (Constructive):} We must find irrational numbers $a$ and $b$ such that $a^b = 2$. Let $a = \sqrt{3}$ and $b = \log_3(4)$. 
    First, we prove $a$ is irrational. It is a well-known theorem that the square root of any prime number is irrational, so $\sqrt{3} \notin \mathbb{Q}$. 
    Second, we prove $b$ is irrational. Assume for contradiction that $\log_3(4) = \frac{p}{q}$ for some $p, q \in \mathbb{Z}^+$ with $q \neq 0$. Then $3^{p/q} = 4 \implies 3^p = 4^q$. The left side is an odd integer (a power of 3), and the right side is an even integer (a power of 4). An odd integer cannot equal an even integer, which is a contradiction. Thus, $b$ is irrational.
    Finally, we evaluate $a^b$:
    \[ a^b = (\sqrt{3})^{\log_3(4)} = (3^{1/2})^{\log_3(4)} = 3^{\frac{1}{2} \log_3(4)} = 3^{\log_3(4^{1/2})} = 3^{\log_3(2)} = 2 \]
    Since $2 \in \mathbb{Q}$, we have constructively proven the existence of such irrational numbers $a$ and $b$. \qedhere

    % Solution 2
    \item \textbf{Proof (Constructive):} Let the six distinct positive integers be $a, b, c, d, e, f$. Let their sum be $S$. The condition that $a$ divides the sum of the remaining integers means $a | (S - a)$, which implies $a | S$. Thus, $S$ must be a multiple of each of the six integers. This implies that $\frac{S}{a}, \frac{S}{b}, \frac{S}{c}, \frac{S}{d}, \frac{S}{e}, \frac{S}{f}$ are distinct positive integers. 
    Summing these integers gives $S = \frac{S}{x_1} + \frac{S}{x_2} + \dots + \frac{S}{x_6}$. Dividing by $S$ yields:
    \[ 1 = \frac{1}{x_1} + \frac{1}{x_2} + \frac{1}{x_3} + \frac{1}{x_4} + \frac{1}{x_5} + \frac{1}{x_6} \]
    We need to find six distinct unit fractions that sum to $1$. We know $1 = \frac{1}{2} + \frac{1}{3} + \frac{1}{6}$. We can expand this using the identity $\frac{1}{n} = \frac{1}{n+1} + \frac{1}{n(n+1)}$.
    Expand $\frac{1}{3} = \frac{1}{4} + \frac{1}{12}$. Now we have $1 = \frac{1}{2} + \frac{1}{4} + \frac{1}{6} + \frac{1}{12}$ (four fractions).
    Expand $\frac{1}{12} = \frac{1}{18} + \frac{1}{36}$. Now we have $1 = \frac{1}{2} + \frac{1}{4} + \frac{1}{6} + \frac{1}{18} + \frac{1}{36}$ (five fractions).
    Expand $\frac{1}{36} = \frac{1}{54} + \frac{1}{108}$. Now we have $1 = \frac{1}{2} + \frac{1}{4} + \frac{1}{6} + \frac{1}{18} + \frac{1}{54} + \frac{1}{108}$ (six fractions).
    Let $S = 108$. The integers are $a = \frac{108}{2} = 54$, $b = \frac{108}{4} = 27$, $c = \frac{108}{6} = 18$, $d = \frac{108}{18} = 6$, $e = \frac{108}{54} = 2$, $f = \frac{108}{108} = 1$. The set $\{1, 2, 6, 18, 27, 54\}$ satisfies the condition. \qedhere

    % Solution 3
    \item \textbf{Proof (Non-Constructive via IVT):} Let $f(x) = x^5 - 3x^2 + 1$. Because $f(x)$ is a polynomial, it is continuous everywhere on $\mathbb{R}$, and specifically continuous on the closed interval $[0, 1]$. 
    We evaluate the function at the endpoints of the interval:
    \[ f(0) = 0^5 - 3(0)^2 + 1 = 1 > 0 \]
    \[ f(1) = 1^5 - 3(1)^2 + 1 = 1 - 3 + 1 = -1 < 0 \]
    Since $f(1) < 0 < f(0)$, the Intermediate Value Theorem guarantees the existence of at least one real number $c \in (0, 1)$ such that $f(c) = 0$. Thus, the equation has a real root in the interval. \qedhere

    % Solution 4
    \item \textbf{Proof (Pigeonhole Principle):} Consider a square of side length $2$. We can partition this square into $4$ smaller, non-overlapping squares, each of side length $1$. 
    We are placing $5$ points (pigeons) into the $4$ smaller squares (holes). By the Pigeonhole Principle, since $5 > 4$, at least one of the smaller squares must contain at least two points.
    The maximum distance between any two points within a square of side length $1$ is the length of its diagonal. By the Pythagorean theorem, the diagonal $d$ is $\sqrt{1^2 + 1^2} = \sqrt{2}$. Therefore, the two points in the same smaller square must be at a distance of at most $\sqrt{2}$ from each other. \qedhere

    % Solution 5
    \item \textbf{Proof by Contradiction:} Assume, to the contrary, that there exists a real number $x$ such that $x^4 + 2x^2 + 2 = 2x$. Rearranging the terms to one side yields:
    \[ x^4 + 2x^2 - 2x + 2 = 0 \]
    We can decompose the terms to form perfect squares:
    \[ x^4 + x^2 + (x^2 - 2x + 1) + 1 = 0 \]
    \[ x^4 + x^2 + (x - 1)^2 + 1 = 0 \]
    For any real number $x$, we know that $x^4 \ge 0$, $x^2 \ge 0$, and $(x - 1)^2 \ge 0$. Therefore, the sum $x^4 + x^2 + (x - 1)^2 + 1$ must be strictly greater than or equal to $1$. 
    \[ x^4 + x^2 + (x - 1)^2 + 1 \ge 1 > 0 \]
    This contradicts the derived equation that the sum equals $0$. Thus, no such real number $x$ exists. \qedhere

    % Solution 6
    \item \textbf{Proof by Contradiction (Parity):} Assume there exists an integer $n$ such that $n^6 + n^3 + 1$ is even. By the definition of parity, $n$ must be either even or odd. We examine both cases:
    \begin{itemize}
        \item \textbf{Case 1 ($n$ is even):} If $n$ is even, then $n^6$ is even and $n^3$ is even. The sum of two even numbers is even, so $n^6 + n^3$ is even. Adding $1$ (an odd number) to an even number results in an odd number. Thus, $n^6 + n^3 + 1$ is odd.
        \item \textbf{Case 2 ($n$ is odd):} If $n$ is odd, then $n^6$ is odd and $n^3$ is odd. The sum of two odd numbers is even, so $n^6 + n^3$ is even. Adding $1$ to an even number results in an odd number. Thus, $n^6 + n^3 + 1$ is odd.
    \end{itemize}
    In all possible cases for $n \in \mathbb{Z}$, the expression evaluates to an odd integer. This contradicts our assumption that it is even. Therefore, no such integer $n$ exists. \qedhere

    % Solution 7
    \item \textbf{Proof by Contradiction (Modular Arithmetic):} Assume there exist odd integers $a, b, c$ such that $a^2 + b^2 + c^2 = 8k + 7$ for some integer $k$. 
    Since $a, b, c$ are odd, they can be written as $a = 2x+1$, $b = 2y+1$, $c = 2z+1$ for some integers $x, y, z$. 
    Squaring $a$ yields: $a^2 = (2x+1)^2 = 4x^2 + 4x + 1 = 4x(x+1) + 1$. 
    Notice that $x(x+1)$ is the product of two consecutive integers, which is always even. Let $x(x+1) = 2m$. Then $a^2 = 4(2m) + 1 = 8m + 1$. Thus, the square of any odd integer is congruent to $1 \pmod 8$.
    Applying this to $a, b,$ and $c$:
    \[ a^2 + b^2 + c^2 \equiv 1 + 1 + 1 \pmod 8 \equiv 3 \pmod 8 \]
    However, the right side of the equation is $8k + 7$, which implies:
    \[ 8k + 7 \equiv 7 \pmod 8 \]
    This forces $3 \equiv 7 \pmod 8$, which is a contradiction. Thus, no such odd integers exist. \qedhere

    % Solution 8
    \item \textbf{Proof (Existence and Uniqueness):} Let $f(x) = e^x + x$. Since $e^x$ and $x$ are continuous and differentiable on $\mathbb{R}$, $f(x)$ is continuous and differentiable on $\mathbb{R}$.
    \textit{Existence:} We evaluate $f(x)$ at $x = -1$ and $x = 0$:
    \[ f(-1) = e^{-1} - 1 = \frac{1}{e} - 1 < 0 \quad \text{(since } e \approx 2.718 > 1 \text{)} \]
    \[ f(0) = e^0 + 0 = 1 > 0 \]
    Since $f(-1) < 0 < f(0)$, the Intermediate Value Theorem guarantees the existence of at least one $c \in (-1, 0)$ such that $f(c) = 0$.
    \textit{Uniqueness:} We examine the first derivative: $f'(x) = e^x + 1$. Since $e^x > 0$ for all $x \in \mathbb{R}$, $f'(x) > 1 > 0$ for all $x \in \mathbb{R}$. Because the derivative is strictly positive, $f(x)$ is strictly monotonically increasing on $\mathbb{R}$. A strictly increasing function can intersect the x-axis at most once. Thus, the root is unique. \qedhere

    % Solution 9
    \item \textbf{Proof (Logical Equivalence):} We start with the left-hand side and apply standard logical equivalences:
    \[ \sim(\exists x \in S, P(x) \implies Q(x)) \]
    By De Morgan's Laws for Quantifiers, negating an existential quantifier results in a universal quantifier:
    \[ \equiv \forall x \in S, \sim(P(x) \implies Q(x)) \]
    By the definition of material implication, $P(x) \implies Q(x)$ is logically equivalent to $\sim P(x) \lor Q(x)$:
    \[ \equiv \forall x \in S, \sim(\sim P(x) \lor Q(x)) \]
    Applying De Morgan's Laws for Propositional Logic to the inner statement:
    \[ \equiv \forall x \in S, (\sim(\sim P(x)) \land \sim Q(x)) \]
    By the law of double negation:
    \[ \equiv \forall x \in S, P(x) \land \sim Q(x) \]
    This matches the right-hand side, establishing the logical equivalence. \qedhere

    % Solution 10
    \item \textbf{Proof (Existence and Uniqueness of Measure):} We must prove $\exists! \mathbb{Q}$ such that $\mathbb{Q}(\omega_1) = q > 0$, $\mathbb{Q}(\omega_2) = 1 - q > 0$, and the discounted stock price is a martingale under $\mathbb{Q}$. The martingale condition is:
    \[ S_0 = \frac{1}{1+r} \mathbb{E}^\mathbb{Q}[S_1] \]
    Substitute the given values ($S_0 = 50, r = 0.1, S_1(\omega_1) = 60, S_1(\omega_2) = 40$):
    \[ 50 = \frac{1}{1 + 0.1} [q(60) + (1-q)(40)] \]
    \[ 50(1.1) = 60q + 40 - 40q \]
    \[ 55 = 20q + 40 \]
    \[ 15 = 20q \implies q = \frac{15}{20} = 0.75 \]
    Since $q = 0.75$, it follows that $1 - q = 0.25$. Both probabilities are strictly between $0$ and $1$, satisfying the definition of an equivalent probability measure. Because the linear equation $15 = 20q$ has exactly one algebraic solution, the measure $\mathbb{Q}$ is unique. \qedhere

    % Solution 11
    \item \textbf{Proof ($\epsilon-\delta$ Definition):} We must show: $\forall \epsilon > 0, \exists \delta > 0$ such that if $0 < |x - 4| < \delta$, then $|(3x - 2) - 10| < \epsilon$.
    Let $\epsilon > 0$ be arbitrary. We manipulate the target inequality to find a suitable $\delta$:
    \[ |(3x - 2) - 10| < \epsilon \implies |3x - 12| < \epsilon \implies 3|x - 4| < \epsilon \implies |x - 4| < \frac{\epsilon}{3} \]
    \textit{Constructive Existence of $\delta$:} Choose $\delta = \frac{\epsilon}{3}$. Since $\epsilon > 0$, it follows that $\delta > 0$. 
    Assume $0 < |x - 4| < \delta$. Then:
    \[ |(3x - 2) - 10| = |3x - 12| = 3|x - 4| < 3\delta = 3\left(\frac{\epsilon}{3}\right) = \epsilon \]
    Thus, the limit exists and equals $10$. \qedhere

    % Solution 12
    \item \textbf{Solution (Flaw Identification):} The flaw lies in the domain of the variable $x$. The prompt explicitly asks to prove the existence of a \textit{real number} $x \in \mathbb{R}$. The proof provides the example $x = i$, which is an imaginary number ($i \notin \mathbb{R}$). Therefore, the proof fails to satisfy the domain constraint of the existential quantifier. \qedhere

    % Solution 13
    \item \textbf{Proof (Constructive):} Let $a = 2$, which is clearly a rational number. Let $b = \log_2(\sqrt{3})$. 
    First, we prove $b$ is irrational. Assume for contradiction that $\log_2(\sqrt{3}) = \frac{p}{q}$ for $p, q \in \mathbb{Z}^+$. Then $2^{p/q} = \sqrt{3} \implies 2^{2p/q} = 3 \implies 2^{2p} = 3^q$. The left side is even, and the right side is odd, which is a contradiction. Thus, $b$ is irrational.
    Now, evaluate $a^b$:
    \[ a^b = 2^{\log_2(\sqrt{3})} = \sqrt{3} \]
    It is a known theorem that $\sqrt{3}$ is irrational. Thus, we have found a rational $a$ and irrational $b$ such that $a^b$ is irrational. \qedhere

    % Solution 14
    \item \textbf{Proof (Constructive):} We must find a prime $p$ such that $p, p+2$, and $p+6$ are all prime. We use heuristic search over small primes.
    Let $p = 5$. 
    The second number is $p + 2 = 5 + 2 = 7$, which is prime.
    The third number is $p + 6 = 5 + 6 = 11$, which is prime.
    Since $5, 7$, and $11$ are all prime numbers, the existence of such a prime $p$ is proven. \qedhere

    % Solution 15
    \item \textbf{Proof (Non-Constructive via IVT):} Let $f(x) = \cos(x) - x^2$. Since $\cos(x)$ and polynomials are continuous on $\mathbb{R}$, $f(x)$ is continuous on $\mathbb{R}$.
    Evaluate $f(x)$ at $x = 0$ and $x = \frac{\pi}{2}$:
    \[ f(0) = \cos(0) - 0^2 = 1 > 0 \]
    \[ f\left(\frac{\pi}{2}\right) = \cos\left(\frac{\pi}{2}\right) - \left(\frac{\pi}{2}\right)^2 = 0 - \frac{\pi^2}{4} < 0 \]
    By the Intermediate Value Theorem, since $f(\frac{\pi}{2}) < 0 < f(0)$, there exists at least one root $c_1 \in (0, \frac{\pi}{2})$.
    Now evaluate $f(x)$ at $x = -\frac{\pi}{2}$:
    \[ f\left(-\frac{\pi}{2}\right) = \cos\left(-\frac{\pi}{2}\right) - \left(-\frac{\pi}{2}\right)^2 = 0 - \frac{\pi^2}{4} < 0 \]
    By the IVT, since $f(-\frac{\pi}{2}) < 0 < f(0)$, there exists at least one root $c_2 \in (-\frac{\pi}{2}, 0)$.
    Since $c_1 > 0$ and $c_2 < 0$, $c_1 \neq c_2$. Thus, there exist at least two distinct real roots. \qedhere

    % Solution 16
    \item \textbf{Proof (Pigeonhole Principle):} By the Division Algorithm, any integer $k$ can be expressed as $k = qn + r$, where $0 \le r \le n-1$. Thus, there are exactly $n$ possible remainders when dividing by $n$.
    We are given a set of $n+1$ integers (the pigeons). The possible remainders modulo $n$ are the $n$ equivalence classes (the holes). 
    Since $n+1 > n$, the Pigeonhole Principle guarantees that at least two distinct integers in our set, say $a$ and $b$, must have the same remainder modulo $n$. 
    Therefore, $a \equiv b \pmod n$, which implies $a - b \equiv 0 \pmod n$. Thus, their difference is a multiple of $n$. \qedhere

    % Solution 17
    \item \textbf{Proof by Contradiction:} Assume there exist real numbers $x$ and $y$ such that $x^2 + y^2 - 4x - 6y + 14 = 0$. We complete the square for both $x$ and $y$:
    \[ (x^2 - 4x + 4) - 4 + (y^2 - 6y + 9) - 9 + 14 = 0 \]
    \[ (x - 2)^2 + (y - 3)^2 - 13 + 14 = 0 \]
    \[ (x - 2)^2 + (y - 3)^2 + 1 = 0 \]
    \[ (x - 2)^2 + (y - 3)^2 = -1 \]
    For any real numbers $x$ and $y$, $(x - 2)^2 \ge 0$ and $(y - 3)^2 \ge 0$. The sum of two non-negative real numbers must be non-negative. However, the equation states their sum is $-1$, which is a contradiction. Thus, no such real numbers exist. \qedhere

    % Solution 18
    \item \textbf{Proof by Contradiction (Parity/Modulo):} Assume there exist integers $a$ and $b$ such that $a^2 - b^2 = 2$. Factoring the difference of squares yields:
    \[ (a - b)(a + b) = 2 \]
    Notice that $(a + b) - (a - b) = 2b$, which is an even integer. This means that $(a - b)$ and $(a + b)$ must have the same parity (both even or both odd). 
    If both are odd, their product is odd, but $2$ is even (contradiction).
    If both are even, then $(a - b) = 2k$ and $(a + b) = 2j$. Their product is $(2k)(2j) = 4kj$. This implies that their product must be a multiple of $4$. However, $2$ is not a multiple of $4$ (contradiction).
    Thus, no such integers $a$ and $b$ exist. \qedhere

    % Solution 19
    \item \textbf{Proof by Contradiction (Modular Arithmetic):} Assume there exist integers $x$ and $y$ such that $x^2 - 3y^2 = 8$. We project this equation modulo $3$:
    \[ x^2 - 3y^2 \equiv 8 \pmod 3 \]
    Since $3y^2 \equiv 0 \pmod 3$ and $8 \equiv 2 \pmod 3$, this simplifies to:
    \[ x^2 \equiv 2 \pmod 3 \]
    We test all possible residues for $x$ modulo $3$:
    If $x \equiv 0 \pmod 3$, $x^2 \equiv 0 \pmod 3$.
    If $x \equiv 1 \pmod 3$, $x^2 \equiv 1 \pmod 3$.
    If $x \equiv 2 \pmod 3$, $x^2 \equiv 4 \equiv 1 \pmod 3$.
    The quadratic residues modulo $3$ are only $0$ and $1$. There is no integer $x$ such that $x^2 \equiv 2 \pmod 3$. This is a contradiction. Thus, no such integers exist. \qedhere

    % Solution 20
    \item \textbf{Proof (Existence and Uniqueness):} 
    \textit{Existence:} Let $A = \emptyset$. For any set $B$, the intersection of the empty set with any set is the empty set. Thus, $\emptyset \cap B = \emptyset$, satisfying the condition $A \cap B = A$.
    \textit{Uniqueness:} Assume there exist two sets $A_1$ and $A_2$ that satisfy the property. Then for all sets $B$, $A_1 \cap B = A_1$ and $A_2 \cap B = A_2$. 
    Since the property holds for all $B$, let $B = A_2$ for the first equation: $A_1 \cap A_2 = A_1$.
    Let $B = A_1$ for the second equation: $A_2 \cap A_1 = A_2$.
    Since set intersection is commutative ($A_1 \cap A_2 = A_2 \cap A_1$), it follows that $A_1 = A_2$. Thus, the set is unique. \qedhere

    % Solution 21
    \item \textbf{Proof (Disproving via Negation):} To prove the statement is false, we must prove its logical negation is true. The negation of $\exists x \in \mathbb{R}, \forall y \in \mathbb{R}, x + y = 0$ is:
    \[ \forall x \in \mathbb{R}, \exists y \in \mathbb{R}, x + y \neq 0 \]
    Let $x \in \mathbb{R}$ be arbitrary. We must constructively find a $y$ such that $x + y \neq 0$. 
    Choose $y = -x + 1$. Since $x \in \mathbb{R}$, $y \in \mathbb{R}$. 
    Substitute $y$ into the expression: $x + y = x + (-x + 1) = 1$. 
    Since $1 \neq 0$, the condition holds. Because $x$ was arbitrary, the negated statement is true, meaning the original statement is false. \qedhere

    % Solution 22
    \item \textbf{Proof by Contradiction (Finance):} Assume there exists an equivalent risk-neutral measure $\mathbb{Q}$ such that $0 < q < 1$, where $q = \mathbb{Q}(\omega_1)$. The martingale condition requires:
    \[ S_0 = \frac{1}{1+r} [q(u \cdot S_0) + (1-q)(d \cdot S_0)] \]
    Divide by $S_0$ (since $S_0 > 0$):
    \[ 1+r = q \cdot u + d - q \cdot d \]
    \[ 1+r - d = q(u - d) \implies q = \frac{1+r-d}{u-d} \]
    We are given that $u = 1+r$. Substitute $u$ for $1+r$ in the numerator:
    \[ q = \frac{u - d}{u - d} = 1 \]
    However, the definition of an equivalent measure requires $0 < q < 1$. Since $q = 1$, it is not strictly less than $1$. This contradicts the assumption that an equivalent risk-neutral measure exists. \qedhere

    % Solution 23
    \item \textbf{Proof ($\epsilon-\delta$ Definition):} We must show: $\forall \epsilon > 0, \exists \delta > 0$ such that if $0 < |x - 2| < \delta$, then $|x^2 - 4| < \epsilon$.
    Let $\epsilon > 0$ be arbitrary. Notice that $|x^2 - 4| = |x - 2||x + 2|$. We need to bound the term $|x + 2|$. 
    Assume $\delta \le 1$. Then $|x - 2| < 1 \implies -1 < x - 2 < 1 \implies 1 < x < 3$. 
    Adding $2$ to all parts yields $3 < x + 2 < 5$. Thus, $|x + 2| < 5$.
    Now, $|x^2 - 4| = |x - 2||x + 2| < 5|x - 2|$. We want this to be less than $\epsilon$, so we need $5|x - 2| < \epsilon \implies |x - 2| < \frac{\epsilon}{5}$.
    \textit{Constructive Existence of $\delta$:} Choose $\delta = \min(1, \frac{\epsilon}{5})$. Since $\epsilon > 0$, $\delta > 0$.
    Assume $0 < |x - 2| < \delta$. Then $|x - 2| < 1$ (which guarantees $|x + 2| < 5$) and $|x - 2| < \frac{\epsilon}{5}$.
    \[ |x^2 - 4| = |x - 2||x + 2| < \left(\frac{\epsilon}{5}\right)(5) = \epsilon \]
    Thus, the limit exists and equals $4$. \qedhere

    % Solution 24
    \item \textbf{Solution (Flaw Identification):} The proof relies on a fundamental logical fallacy regarding Proof by Contradiction. In a valid proof by contradiction, one assumes the negation of the statement and derives a \textit{false} statement (a contradiction), which implies the negation is false and the original statement is true. 
    In this flawed proof, the author assumes the negation and derives a \textit{true} statement ($0 \le 1$). Deriving a true statement from an assumption does not invalidate the assumption. In fact, it provides no logical grounds to reject the negation. The original statement ($\exists x, x > x+1$) is definitively false. \qedhere

    % Solution 25
    \item \textbf{Proof (Constructive):} We must show that the set $S = \{q \in \mathbb{Q} : 0 < q < 1\}$ is infinite. 
    Consider the sequence of numbers defined by $q_n = \frac{1}{n}$ for integers $n \ge 2$. 
    For any integer $n \ge 2$, $n > 1 \implies \frac{1}{n} < 1$. Also, since $n > 0$, $\frac{1}{n} > 0$. Thus, $0 < q_n < 1$.
    Furthermore, $q_n$ is the ratio of two integers, so $q_n \in \mathbb{Q}$. 
    Since there are infinitely many integers $n \ge 2$, and each produces a distinct fraction $\frac{1}{n}$, there exist infinitely many rational numbers in the interval $(0, 1)$. \qedhere

    % Solution 26
    \item \textbf{Proof (Constructive):} We must find a positive integer $n$ such that $n! + 1 = k^2$ for some integer $k$. We test small values of $n$:
    If $n = 1$, $1! + 1 = 2$ (not a square).
    If $n = 2$, $2! + 1 = 3$ (not a square).
    If $n = 3$, $3! + 1 = 7$ (not a square).
    If $n = 4$, $4! + 1 = 24 + 1 = 25$. 
    Notice that $25 = 5^2$, which is a perfect square. Thus, $n = 4$ satisfies the condition, proving existence. \qedhere

    % Solution 27
    \item \textbf{Proof (Non-Constructive via IVT):} Let $g(x) = f(x) - x$. Since $f(x)$ is continuous on $[a, b]$ and $x$ is continuous, $g(x)$ is continuous on $[a, b]$.
    We evaluate $g(x)$ at the endpoints $a$ and $b$:
    \[ g(a) = f(a) - a \]
    Since the codomain of $f$ is $[a, b]$, we know $f(a) \ge a$. Thus, $g(a) \ge 0$.
    \[ g(b) = f(b) - b \]
    Similarly, since the codomain is $[a, b]$, $f(b) \le b$. Thus, $g(b) \le 0$.
    If $g(a) = 0$, then $f(a) = a$, and $c = a$ is a fixed point.
    If $g(b) = 0$, then $f(b) = b$, and $c = b$ is a fixed point.
    If $g(a) > 0$ and $g(b) < 0$, then $g(b) < 0 < g(a)$. By the Intermediate Value Theorem, there exists at least one $c \in (a, b)$ such that $g(c) = 0$. 
    Since $g(c) = f(c) - c = 0 \implies f(c) = c$, a fixed point must exist. \qedhere

    % Solution 28
    \item \textbf{Proof (Generalized Pigeonhole Principle):} We have $n = 101$ items (pigeons) and $m = 10$ containers (holes). The Generalized Pigeonhole Principle states that if $n$ items are placed into $m$ containers, at least one container must hold at least $\lceil \frac{n}{m} \rceil$ items.
    Evaluating the ceiling function:
    \[ \left\lceil \frac{101}{10} \right\rceil = \lceil 10.1 \rceil = 11 \]
    Therefore, at least one container must contain $11$ or more items. \qedhere

    % Solution 29
    \item \textbf{Proof by Contradiction:} Assume there exist strictly positive integers $x$ and $y$ such that $x^2 - y^2 = 1$. Factoring the left side yields:
    \[ (x - y)(x + y) = 1 \]
    Since $x$ and $y$ are strictly positive integers ($x \ge 1, y \ge 1$), their sum must be $x + y \ge 2$. 
    The only integer factors of $1$ are $(1)(1)$ and $(-1)(-1)$. Thus, we must have $x + y = 1$ or $x + y = -1$. 
    However, we already established that $x + y \ge 2$. This is a contradiction ($2 \le x+y = 1$). Thus, no such positive integers exist. \qedhere

    % Solution 30
    \item \textbf{Proof (Pigeonhole Principle):} Let $S$ be a set of $5$ distinct positive integers. The number of distinct non-empty subsets of $S$ is $2^5 - 1 = 32 - 1 = 31$. These $31$ subsets are our "pigeons".
    For each subset, we calculate the sum of its elements modulo $30$. The possible equivalence classes modulo $30$ are $\{0, 1, 2, \dots, 29\}$. There are exactly $30$ possible sums modulo $30$. These are our "holes".
    Since we have $31$ subsets and only $30$ possible sums modulo $30$, the Pigeonhole Principle ($31 > 30$) guarantees that at least two distinct non-empty subsets, say $A$ and $B$, must map to the same sum modulo $30$. 
    Therefore, the sum of the elements in $A$ is congruent to the sum of the elements in $B$ modulo $30$. \qedhere

\end{enumerate}

\clearpage

\section{Concept Checklist \& Question Mapping (Reference)}
\label{sec:checklist_end}
\begin{itemize}
    \item \textbf{C1: Constructive Existence (Algebraic/Rationality)} -- Problems 1, 13, 25
    \item \textbf{C2: Constructive Existence (Primes/Divisibility/Factorials)} -- Problems 2, 14, 26
    \item \textbf{C3: Non-Constructive Existence (Intermediate Value Theorem)} -- Problems 3, 8, 15, 27
    \item \textbf{C4: Combinatorial Existence (Pigeonhole Principle)} -- Problems 4, 16, 28, 30
    \item \textbf{C5: Disproving Existence (Algebra \& Sum of Squares)} -- Problems 5, 17, 29
    \item \textbf{C6: Disproving Existence (Parity \& Consecutive Integers)} -- Problems 6, 18
    \item \textbf{C7: Disproving Existence (Modular Arithmetic \& Finite Fields)} -- Problems 7, 19
    \item \textbf{C8: Uniqueness Proofs ($\exists!$)} -- Problems 8, 10, 20
    \item \textbf{C9: Set-Theoretic \& Logical Equivalences (De Morgan's)} -- Problems 9, 21
    \item \textbf{C10: Abstract Real-World Modeling (Quant Finance/Risk-Neutral Measures)} -- Problems 10, 22
    \item \textbf{C11: Formal Limit Existence ($\epsilon-\delta$ Nested Quantifiers)} -- Problems 11, 23
    \item \textbf{C12: Advanced Diagnostic Testing (Finding the Flaw)} -- Problems 12, 24
\end{itemize}

\end{document}